\begin{small}
\begin{longtable}{|>{\raggedright\arraybackslash}p{0.18\textwidth}|>{\raggedright\arraybackslash}p{0.4\textwidth}|>{\raggedright\arraybackslash}p{0.4\textwidth}|}
\caption{Consolidated Implementation Blueprint Summary}\label{tab:blueprint_summary} \\
\hline
\multicolumn{3}{|c|}{\textbf{Consolidated Implementation Blueprint Summary}} \\
\hline
\textbf{Component} & \textbf{Definitive Specification} & \textbf{Justification} \\
\hline
\endfirsthead
\hline
\textbf{Component} & \textbf{Definitive Specification} & \textbf{Justification} \\
\hline
\endhead
% Table rows start here
\textbf{Unit System} & A non-dimensional system is used where the gravitational constant $G=1$. \newline \textbullet{} \textbf{Characteristic Time:} $T=t_{max}$ (max time in dataset). \newline \textbullet{} \textbf{Characteristic Length:} $D=\sqrt[3]{GT^{2}}$. & Simplifies the governing physical equations and improves the numerical stability of the model during training. \\
\hline
\textbf{Backbone} & \textbf{SQR-SkipResNet} with \textbf{12 hidden layers} and a \textbf{width of 256 neurons}. \newline \textbullet{} \textbf{Odd Layers:} $\mathbf{x}^{(l)} = f_{\text{sr}}^{(l)}(\mathbf{x}^{(l-1)}) + (\mathrm{LayerNorm}(\mathbf{x}^{(l-1)}))^2$. \newline \textbullet{} \textbf{Even Layers:} Plain block, $\mathbf{x}^{(l)} = f_{\text{pn}}^{(l)}(\mathbf{x}^{(l-1)})$. \newline \textbullet{} \textbf{Layer Norm:} Applied to skip path on odd layers only, with \texttt{elementwise\_affine=True} and \texttt{eps=1e-5}. & The deep architecture provides the expressive power needed for chaotic dynamics. The SQR-SkipResNet variant with LayerNorm provides superior gradient stability for non-smooth dynamics, which is ideal for modeling close encounters. \\
\hline
\textbf{Activation} & \textbf{SIREN ($\sin(x)$)} for all hidden layers. & The sine function is infinitely differentiable ($C^{\infty}$), which is a mathematical necessity for accurately computing the second-order derivatives ($\frac{d^{2}r}{dt^{2}}$) required in the physics loss term. \\
\hline
\textbf{Initialization} & A principled, two-part SIREN initialization scheme: \newline \textbullet{} \textbf{First Layer:} \texttt{first\_omega\_0 = 30}. Weights from $U[-\frac{1}{n_{in}}, \frac{1}{n_{in}}]$. \newline \textbullet{} \textbf{Hidden Layers:} \texttt{hidden\_omega\_0 = 1}. Weights from $U[-\frac{\sqrt{6/n_{in}}}{1}, \frac{\sqrt{6/n_{in}}}{1}]$. & A high initial frequency (\texttt{first\_omega\_0}) helps overcome spectral bias. The specific scaling for hidden layers ensures a stable distribution of activations, making the deep network trainable. \\
\hline
\textbf{Input Vector} & A \textbf{14-dimensional vector} \texttt{[t, IC]}, where \texttt{t} is time normalized to $[-1,1]$ and \texttt{IC} is a 13-component vector of initial conditions. The IC vector contains: \newline \textbullet{} 3 masses. \newline \textbullet{} 5 physical invariants. \newline \textbullet{} 5 values for 3 trigonometrically encoded angles. & This formulation provides a physically-grounded and injective (unambiguous) representation of the initial state, which is critical for a conditional network to learn correctly. \\
\hline
\textbf{Output} & A \textbf{9-dimensional position vector} \texttt{[r1(t), r2(t), r3(t)]} with a linear (no activation) final layer. Velocities ($v=\frac{dr}{dt}$) and accelerations ($a=\frac{d^{2}r}{dt^{2}}$) are derived via \textbf{Automatic Differentiation (AD)}. & Predicting only positions simplifies the network architecture. Deriving kinematics via AD hard-codes their fundamental physical relationship, ensuring they are satisfied by construction to machine precision. \\
\hline
\textbf{Data Loss ($\mathcal{L}_{data}$)} & \textbf{Mean Absolute Error (MAE / L1 Norm)} of the predicted positions against the ground truth data. & MAE is more robust to the occasional large deviations and outliers that are common in chaotic N-body simulations compared to MSE. \\
\hline
\textbf{Physics Loss ($\mathcal{L}_{physics}$)} & \textbf{Mean Squared Error (MSE)} of the physics residual. The residual is the difference between the network's AD-computed acceleration and the acceleration calculated from Newton's law of gravitation (in $G=1$ units). & This loss enforces the governing ordinary differential equations across the entire solution domain, not just at data points, allowing the model to generalize and make physically consistent predictions. \\
\hline
\textbf{Energy Loss ($\mathcal{L}_{energy}$)} & \textbf{Huber Loss ($\delta=1.0$)} on the deviation between the network's predicted total energy (Hamiltonian) and the pre-computed, constant initial energy ($H_0$), all in $G=1$ units. & Enforces the global conservation of energy, which combats long-term numerical drift and improves the physical plausibility of the simulation over long time horizons. \\
\hline
\end{longtable}
\end{small}
